This phase analyzes the potential consequences of the confirmed
vulnerabilities when considered individually and in combination.
The analysis focuses on demonstrated security impact as well as
realistic attack-chain scenarios that could occur if the
vulnerabilities were present in a production deployment handling
real user data.

\subsection*{Objective}

Phase 4 confirmed the identified vulnerabilities through individual
manual verification and proof-of-concept testing. Phase 5 considers
how these confirmed weaknesses could interact with one another and
how their combined impact could affect the application's
confidentiality, integrity, and access-control mechanisms.

The attack-chain scenarios presented below represent potential
combined impacts based on the confirmed vulnerabilities. They are
not presented as additional exploitation activities unless
specifically demonstrated during the assessment.

\subsection*{Attack Chain Scenarios}

\subsubsection*{Scenario 1: Unauthenticated Database Access via SQL Injection}

The SQL Injection vulnerability in the product search endpoint
(SQLI-02) does not require authentication. During testing, SQL
Injection was used to identify the SQLite database structure and
retrieve records from the \texttt{Users} table, including email
addresses and password hashes.

In a production environment, unauthorized access to user records
could expose sensitive account information. If users reused weak
passwords or if the exposed password hashes were successfully
cracked, the information could potentially be used in credential-
stuffing attacks against other services.

The assessment demonstrated unauthorized database access through
SQLI-02; credential cracking and attacks against external services
were not performed as part of this assessment.

\subsubsection*{Scenario 2: Privilege Escalation to Administrative Access}

The Broken Access Control finding (BAC) allows an unauthenticated
visitor to register an account with the \texttt{admin} role. During
testing, the newly created account successfully accessed the
Administration page.

SQLI-01 provides a separate authentication-bypass path that was also
demonstrated during testing. Therefore, the application contains
multiple independent weaknesses that can result in unauthorized
administrative access.

Once administrative access is obtained, the application's
administration functionality may provide access to user data,
product information, and other administration-only functionality.
The demonstrated administrative access represents a significant
failure of the application's authorization controls.

\subsubsection*{Scenario 3: Client-Side Impact via Reflected XSS and CSRF}

The Reflected XSS finding demonstrates that attacker-controlled
JavaScript can execute in a victim's browser through a crafted
search request.

The CSRF finding separately demonstrates that a cross-site request
can cause a profile update while the victim is authenticated.

In combination, these weaknesses could increase the impact of a
client-side attack by allowing attacker-controlled script execution
and unauthorized state-changing requests under suitable conditions.

The assessment confirmed the XSS and CSRF vulnerabilities
individually. A complete combined XSS-to-CSRF attack chain, including
session-data extraction or other browser-side credential theft, was
not separately demonstrated and is therefore considered a potential
impact scenario rather than a confirmed exploitation result.

\subsubsection*{Scenario 4: Account Takeover via Weak Password Recovery}

The weak password recovery mechanism (AUTH) allows an attacker who
can obtain, research, or correctly guess the relevant security-
question answer to reset a target account's password.

During testing, the password recovery weakness was successfully
demonstrated against the tested account. If the same weakness were
present for an account with administrative privileges, successful
exploitation could potentially provide unauthorized administrative
access.

This represents an additional path to account compromise that is
independent of the BAC and SQLI-01 findings.

\subsubsection*{Scenario 5: Unauthorized Data Exposure via IDOR}

The IDOR finding allows an authenticated user to view another user's
basket contents by manipulating the basket identifier.

The assessment demonstrated unauthorized access to another user's
basket data. In a production environment, automated enumeration of
predictable or accessible identifiers could potentially increase the
number of affected accounts and expose additional shopping-related
information.

The assessment did not perform large-scale enumeration of basket
identifiers; therefore, such broader data harvesting is considered a
potential impact scenario.

\subsection*{Consolidated Business Impact}

If the confirmed vulnerabilities existed in a production application
handling real customer data, the combined impact could include:

\begin{itemize}

    \item \textbf{Unauthorized administrative access} through
    multiple independent weaknesses, including the demonstrated BAC
    and SQLI-01 attack paths and the potential impact of the AUTH
    weakness against privileged accounts.

    \item \textbf{Unauthorized access to user database records}
    through SQLI-02, including email addresses and password hashes
    retrieved during testing.

    \item \textbf{Client-side attacks against legitimate users}
    through Reflected XSS, with additional unauthorized state changes
    potentially possible when combined with the confirmed CSRF
    weakness.

    \item \textbf{Exposure of individual users' shopping and basket
    data} through IDOR.

    \item \textbf{Erosion of user trust} through Open Redirect, which
    could potentially be used to support phishing campaigns by
    redirecting users from the legitimate application domain to an
    attacker-controlled destination.

\end{itemize}

\subsection*{Conclusion}

The individual findings documented in Phase 4 are significant on
their own, and several provide overlapping paths to high-value
outcomes such as unauthorized administrative access or sensitive
data exposure.

The assessment demonstrated multiple independent weaknesses affecting
authentication and authorization controls. This indicates that
remediation should address the underlying security controls rather
than relying on a single isolated fix. For example, addressing
SQLI-01 alone would not remove the administrative-access risk created
by the BAC finding, while the AUTH weakness represents an additional
potential path to account compromise.

Detailed remediation guidance for each individual finding, along with
consolidated recommendations, is provided in
Chapter~\ref{chap:phase7}.